%% ==================== 模型评价 模板 (v1.5.6, BZD 数模社 2026 规范) ====================
%%
%% 借鉴 BZD 数模社 第六章 模板: 总结 + 优点 + 不足
%% 注意: 4 子节结构 (BZD 常见形式 A) 或 3 子节 (形式 B)
%%
%% BZD 原则:
%% - 优点: "具体做法 + 形成的优势 + 证据或效果" (避免 "精度高、创新性强" 套话)
%% - 不足: "不足来源 + 可能影响 + 影响范围" (避免足以否定全文的致命问题)
%% - 改进: 必须逐项对应不足
%% - 推广: 说明对象 + 变量替换 + 参数标定 + 约束调整 + 所需数据
%%
%% 跑题用户必须:
%% 1. 删掉所有 % 注释
%% 2. 替换【】占位符为自己的实际内容
%% 3. 删 下面的 \EvalSlot{...} 行 (check 15 抓这个 inline 【】)

\section{模型评价、改进与推广}

\subsection{模型总结}

本文针对【研究问题】, 分别建立【模型 1】、【模型 2】、【模型 3】和【模型 4】。各模型依次完成【问题一目标】、【问题二目标】、【问题三目标】和【问题四目标】, 并通过【检验方法: 5 折交叉验证 + 灵敏度分析 + 稳健性检验】验证结果可靠性。

问题一通过【方法, 如 多元线性回归 + 随机森林回归】识别【主要因素】; 问题二在问题一的基础上建立【如 ROC 曲线 + 贝叶斯决策】, 推荐【最佳时点】; 问题三在问题二的基础上按【分组标准】分层, 揭示【分组差异】; 问题四在问题三的基础上建立【如 随机森林 + Isolation Forest】, 形成【决策规则 + 比例】。

经检验, 各模型在【数据范围】内具有较好的【准确性 ($R^2 > 0.XX$)、稳定性 (CV 标准差 < 0.XX) 和泛化能力 (测试集误差 < 训练集)】, 可应用于【具体场景或推广对象】。

\subsection{模型的优点}

（1）\textbf{模型具有较强针对性}。针对【题目特点】, 在【经典模型】基础上增加【新变量/新约束/新参数】, 提高了模型与实际问题的匹配程度。【如: 在标准 Logistic 回归基础上引入 BMI $\times$ 孕周交互项, 揭示了 BMI 对 NIPT 检测时点的非线性影响】

（2）\textbf{数据处理较为完整}。针对【数据问题】采用【处理方法: 缺失值 + 异常值 + 标准化】, 降低了数据质量对后续模型的影响。【如: 经清洗后保留 N 个有效样本, 处理后的数据通过了 K-S 正态性检验和方差齐性检验】

（3）\textbf{问题之间联动清晰}。将问题一的【输出: 主要特征/权重/参数】作为问题二的【输入】, 进一步进入问题三的【目标函数或约束】, 最后由问题四形成【综合决策】, 实现了数据流、参数流和结论流的完整串联。【如: 特征重要性 (0.42) 作为问题二的关键特征, 预测值作为问题三的输入参数, 决策概率作为问题四的判据】

（4）\textbf{模型结果可靠}。经【检验方法: 5 折交叉验证 + 灵敏度分析 + 稳健性检验】验证, 【量化指标, 如: $R^2 = 0.XX$, $\text{AUC} = 0.XX$】达到【具体数值】, 主要结论与【现实规律或临床观察】一致。

（5）\textbf{具有推广价值}。调整【参数或约束】后, 可用于【相似对象或场景】, 详见 7.4 节。

\subsection{模型的不足}

（1）受【数据来源: 单一医院 / 单一地区 / 时间跨度短】限制, 未能充分考虑【其他因素】, 可能对【具体结果】产生一定影响。【如: 本文数据来自单家三甲医院, 缺乏多中心验证, 推广到基层医院或不同地区时可能存在偏差】

（2）为保证模型可解, 对【现实过程】进行了【具体简化】, 使模型在【特定情景】下的适用性受到限制。【如: 假设各 BMI 组内部完全独立, 但实际可能存在组间空间相关或遗传背景相似】

（3）模型结果对参数【参数名】较为敏感, 该参数估计误差可能影响最终结论。【如: 风险权重 $w_1:w_2 = 2:1$ 基于临床经验, 不同权重设置下最优方案可能略有变化, 但整体结论 (最佳时点) 保持稳定】

（4）【算法名称】存在一定随机性, 虽然多次独立运行结果较稳定, 但不能从理论上保证全局最优。【如: 差分进化算法可能陷入局部最优, 可通过增加种群规模或迭代次数改善】

\subsection{模型的改进与推广}

\textbf{（1）模型改进方向} \quad 针对上述不足, 本文提出以下改进方向:
\begin{itemize}
  \item 进一步收集【数据类型, 如 多中心数据 / 更长时间序列】, 扩大时间和空间范围, 提高参数估计可靠性
  \item 增加【新变量/新约束】, 进一步描述当前未充分考虑的现实因素
  \item 引入【动态模型 / 鲁棒优化 / 改进算法】, 降低不确定性对结果的影响
  \item 使用【更精细的随机性控制或理论保证】, 解决算法的随机性问题
\end{itemize}

\textbf{（2）模型推广} \quad 本文模型的核心结构是根据【输入变量】建立【目标变量】与【约束】之间的关系, 因此可以推广至【其他对象, 如 类似疾病筛查 / 类似临床决策问题】。推广时需要:
\begin{itemize}
  \item 将【原变量, 如 Y 染色体浓度】替换为【新变量, 如 其他疾病相关游离 DNA 浓度】
  \item 重新标定【关键参数, 如 BMI 阈值、风险权重】
  \item 增加【新约束, 如 不同疾病的阈值和决策规则】
  \item 补充【新数据, 如 目标疾病的检测数据】
\end{itemize}

推广后的模型可用于解决【类似任务, 如 其他染色体非整倍体筛查 / 其他无创产前检测】, 具有较好的实际应用价值。

% 跑题后必须删下面这一行 (check 15 占位符自检)
\textbf{【这里粘贴完整的 4 子节内容, 不超过 1-2 页, 删掉这一行】}
